\title{Direct Preference Optimization: Your Language Model is Secretly a Reward Model}

\begin{abstract}
Large-scale unsupervised language models (LMs) acquire extensive world knowledge and some reasoning abilities during training, yet precisely directing their output remains challenging due to the fully unsupervised learning paradigm. Methods for improving model steerability typically involve gathering human-annotated comparisons of output quality and subsequently adapting the LM to these judgments, commonly via reinforcement learning from human feedback (RLHF). RLHF, however, introduces significant complexity and instability: it requires first training a reward model that represents human preferences, then optimizing the original unsupervised LM through reinforcement learning to enhance this learned reward signal, all while constraining deviation from the base model. In this work, we propose a novel reward model parameterization for RLHF that allows the derivation of the optimal policy in closed form. This reformulation permits the resolution of the conventional RLHF objective using a straightforward classification loss. We term the resultant method \textit{Direct Preference Optimization} (DPO), which offers robust performance, computational efficiency, and stability. DPO obviates the need for sampling during fine-tuning and significantly reduces the demands of hyperparameter tuning. Experimental results demonstrate that DPO is capable of adapting LMs to align with human preferences on par with, or surpassing, leading existing techniques. Particularly, DPO outperforms PPO-based RLHF approaches for steering the sentiment of generated text, and achieves equal or superior results in summarization and single-turn dialogue tasks, all while markedly simplifying implementation and training.
\end{abstract}

\section{Introduction}
Large-scale unsupervised language models (LMs), when exposed to extensive datasets, demonstrate remarkable and unexpected competencies~\citep{chowdhery2022palm, brown2020language, touvron2023llama, bubeck2023sparks}. Nonetheless, because these models are trained on data produced by humans who possess diverse objectives, interests, and proficiencies, not every learned trait or skill is appropriate for replication. For instance, while it is beneficial for an AI-based programming assistant to \textit{recognize} common coding errors in order to remedy them, it is preferable that the system emulates the (perhaps infrequently observed) expertise found among highly skilled programmers during code generation. Likewise, although a language model should be \textit{informed} of prevalent misconceptions—such as those endorsed by half of a population—we do not intend for it to propagate such errors at a comparable frequency in its outputs. Therefore, distinguishing and promoting the model's \emph{intended outputs and conduct} from its broad array of \textit{internalized knowledge and skills} is imperative for developing AI that is reliable, effective, and controllable~\citep{ouyang2022training}. Present methodologies frequently leverage reinforcement learning (RL) to align LM behavior with human preferences; here, we argue that the RL-based optimization objective traditionally employed in such methods can, in fact, be formulated as a straightforward binary cross-entropy loss, thus streamlining the overall process of preference learning.

In general, current techniques infuse language models with preferred behaviors by utilizing thoughtfully assembled datasets that capture human preferences for safe and useful actions. This preference acquisition follows a primary stage of extensive unsupervised pre-training on large corpora. While supervised fine-tuning—based on human-labeled demonstrations of ideal responses—provides a simple pathway to preference alignment, the leading paradigm involves reinforcement learning from human (or artificial) feedback (RLHF/RLAIF; \citep{christiano2017deep,bai2022constitutional}). Within RLHF, a reward model is trained on annotated preference data, and RL is subsequently used to adjust the language model’s policy to favor outputs assigned high rewards, while constraining deviations from the original pre-trained policy. Although RLHF methods yield models capable of advanced conversation and programming, the RLHF approach is substantially more intricate than direct supervised learning, requiring the training of several LMs, iterative policy sampling, and as a consequence, incurs considerable computational overhead.

In this work, we introduce a method for directly aligning a language model with human preferences, circumventing the need for explicit reward models or reinforcement learning approaches. We present \textit{{Direct Preference Optimization} (DPO)}, an algorithm that, while implicitly targeting the same objective as conventional RLHF methods (i.e., reward maximization under a KL-divergence penalty), is both easy to implement and efficient to train. Conceptually, DPO updates serve to boost the relative log-probabilities of preferred responses over non-preferred ones, incorporating a dynamically computed, example-specific importance weighting. This mechanism guards against the model degradation seen with naïve probability ratio-based objectives. As with prior approaches, DPO utilizes a theoretical preference model (for instance, the Bradley-Terry model; \cite{bradley1952rankanalysis}) to quantify the congruence between a reward function and observed human preference data. Distinctly, whereas established techniques employ the preference model to construct a loss for training a reward model—subsequently guiding a policy to maximize this learned reward—DPO employs a change of variables to directly define the preference loss in terms of the policy. Leveraging datasets of human-labeled preference judgments, DPO trains the policy via a binary cross-entropy objective, yielding an optimal policy with respect to an implicit reward function that closely fits the preference annotations.

Our principal contribution is the development of Direct Preference Optimization (DPO), an accessible algorithm for learning from preferences without relying on reinforcement learning frameworks. Experimental results indicate that DPO matches or outperforms leading alternatives—including RLHF based on PPO—for preference-driven tasks like sentiment modulation, summarization, and dialogue, using language models as large as 6B parameters.

\section{Related Work}

Self-supervised language models, when scaled up, have demonstrated the ability to solve certain tasks in a zero-shot manner \citep{radford2019language} or given only a handful of examples as prompts \citep{gpt3,megatron,chowdhery2022palm}. Nonetheless, their effectiveness on specific applications and the degree to which their outputs align with user intentions are substantially enhanced through fine-tuning with curated collections of instructions paired with human-generated responses \citep{mishra-etal-2022-cross,sanh2022multitask,chung2022scaling,thoppilan2022lamda}. This process, commonly referred to as `instruction-tuning,' extends the generalization capabilities of LLMs beyond the explicit set of instructions included during training and improves their general utility \citep{chung2022scaling}. In contrast to instruction tuning, it is often more feasible to gather \textit{relative} human assessments comparing response quality than to collect extensive expert-labeled demonstrations. Building on this, subsequent research has refined LLMs using preference datasets, thereby advancing model performance in domains such as translation \citep{kreutzer-etal-2018-reliability}, summarization \citep{stiennon2022learning,ziegler2020finetuning}, narrative generation \citep{ziegler2020finetuning}, and following user instructions \citep{ouyang2022training,ramamurthy2023is}. Typically, these approaches involve initially training a reward model to be consistent with observed preference data—often adopting frameworks like the Bradley-Terry model \citep{bradley1952rankanalysis}—and then fine-tuning the language model itself through reinforcement learning to maximize these learned rewards. Reinforcement learning algorithms such as REINFORCE \citep{williams1992reinforce}, proximal policy optimization (PPO; \cite{schulman2017proximal}), or their derivatives are frequently employed \citep{ramamurthy2023is}. Parallel work uses LLMs already instruction-tuned with human feedback as generators for synthetic preference examples that target specific properties like safety or non-harmfulness \citep{bai2022constitutional}, requiring only loosely-specified guidance from humans in the form of written rubrics. Together, these research directions combine efforts on reinforcement learning-based training of language models for diverse objectives~\citep{Ranzato2015SequenceLT,paulus2018a,wu2018learning} and general-purpose methods for harnessing human preferences in machine learning \citep{christiano2017deep,kupcsik2018learning}. While the use of relative human judgments is attractive, the practical implementation of reinforcement learning for fine-tuning large language models poses considerable difficulties; the present work addresses this gap by offering a theoretically principled alternative for optimizing with relative preferences without employing reinforcement learning.

Outside of linguistic applications, the study of learning policies based on preferences has been explored within both bandit frameworks and reinforcement learning paradigms, yielding a variety of proposed methods. In the contextual bandit setting, when the feedback comprises preferences or action rankings instead of numerical rewards, this framework is termed the contextual dueling bandit (CDB; \cite{yue2012karmed,dudik2015contextual}). When absolute rewards are unavailable, theoretical analyses in the CDB literature introduce the concept of a \textit{von Neumann winner}—defined as a policy whose average probability of outperforming \textit{every} other policy is at least 50\% \citep{dudik2015contextual}. Importantly, CDB problems typically assume online provision of preference information, whereas human preference learning usually relies on a predetermined, fixed dataset of preference-annotated action pairs collected offline \citep{yan2022human}. In a similar vein, \textit{preference-based RL} (PbRL) investigates learning via binary feedback generated by an opaque ‘scoring’ function, as opposed to reward signals \citep{BusaFekete2014,ruiz2023dueling}. Many algorithms for PbRL have been introduced, some enabling the reuse of off-policy preference data; yet these generally require the explicit estimation of a latent scoring—or reward—function before the subsequent policy optimization \citep{jain2013learning,BusaFekete2014,christiano2017deep,sadigh2017active,kupcsik2018learning}. In contrast, we introduce a unified approach for policy learning that directly seeks to align a policy with the given preferences.

\section{Preliminaries}\label{section:prelims}

We summarize the RLHF pipeline outlined by \citeauthor{ziegler2020finetuning} (and further developed in \citep{stiennon2022learning, bai2022training, ouyang2022training}), which is generally structured in three stages: 1) supervised fine-tuning (SFT); 2) collection of preference data and reward model learning; and 3) reinforcement learning-based optimization.

\textbf{SFT}: The RLHF process typically initiates by adapting a pre-trained language model to the downstream application (such as dialogue or summarization tasks) through supervised learning using curated, high-quality data. This yields an initial policy denoted as $\pi^\text{SFT}$.

\textbf{Reward Modelling Phase}: In the next step, the supervised fine-tuned model is presented with inputs $x$ and generates pairs of outputs $(y_1, y_2)\sim \pi^\text{SFT}(y \mid x)$. These output pairs are then evaluated by human annotators, who indicate a preference for one completion over the other, recorded as $y_w\succ y_l \mid x$, with $y_w$ representing the favored output and $y_l$ the less preferred option within $(y_1, y_2)$. The assumption is that such preferences are determined by an inaccessible, underlying reward function $r^*(y, x)$. Several models have been proposed for capturing these preferences, with the Bradley-Terry (BT) model \cite{bradley1952rankanalysis} being among the most widely used. Additionally, more general ranking frameworks, such as the Plackett-Luce models \citep{plackett1975analysis, luce2012individual}, can also be applied, especially when multiple ranked completions are provided. Under the BT model, the probability distribution representing human preferences, denoted as $p^*$, is defined as follows:

\begin{equation}\label{eq:bradley-terry}
    p^*(y_1\succ y_2 \mid x)=\frac{\exp\left(r^*(x, y_1)\right)}{\exp\left(r^*(x, y_1)\right) + \exp\left(r^*(x, y_2)\right)}.
\end{equation}

Given a fixed dataset of pairwise preferences $\mathcal{D}=\bigl\{x^{(i)}, y_w^{(i)}, y_l^{(i)}\bigr\}_{i=1}^N$ drawn from $p^*$, one can introduce a parameterized reward function $r_{\phi}(x, y)$ and learn its parameters through maximum likelihood estimation. When this is framed as a binary classification problem, the negative log-likelihood loss is given by:

\begin{equation}\label{eq:reward_model}
    \mathcal{L}_R(r_{\phi}, \mathcal{D}) = -\mathbb{E}_{(x, y_w, y_l)\sim \mathcal{D}}\bigl[\log \sigma(r_{\phi}(x, y_w)- r_{\phi}(x, y_l))\bigr]
\end{equation}

with $\sigma$ denoting the logistic sigmoid. In large language models, the architecture of $r_{\phi}(x, y)$ typically builds on top of the base SFT model $\pi^\text{SFT}(y \mid x)$ by appending a linear projection to the final transformer layer, thereby producing a single scalar reward prediction \cite{ziegler2020finetuning}. To promote lower-variance reward functions, standard practice is to normalize rewards so that $\mathbb{E}_{x,y\sim \mathcal{D}}\left[r_\phi(x, y)\right] = 0$ holds across all $x$.

\textbf{RL Fine-Tuning Phase}: In the subsequent RL stage, the trained reward function acts as the supervisory signal for optimizing the language model policy. Following prior works~\citep{jaques2017sequence, jaques2020human}, the learning objective is formalized as:

\begin{equation}\label{eq:RL}
\max_{\pi_{\theta}}  \mathbb{E}_{x\sim \mathcal{D}, y\sim \pi_{\theta}(y \mid x)}\bigl[r_{\phi}(x, y)\bigr] - \beta\mathbb{D}_{\textrm{KL}}\bigl[\pi_{\theta}(y\mid x)\mid \mid \pi_\text{ref}(y\mid x)\bigr],
\end{equation}

where $\beta$ regulates divergence from the baseline policy $\pi_\text{ref}$, which is chosen as the initial SFT model $\pi^\text{SFT}$. Typically, $\pi_\theta$ is also initialized with $\pi^\text{SFT}$. This regularization is crucial for ensuring that the policy does not stray far from regions where the reward function remains reliable, while also helping to sustain generation diversity and mitigating against mode collapse to a small set of high-reward outputs. Because the language output space is discrete, direct gradient-based optimization of this objective is not possible; thus, reinforcement learning methods are used instead. The prevailing approach \citep{ziegler2020finetuning, stiennon2022learning, bai2022training, ouyang2022training} constructs the following composite reward: ${r(x, y) = r_{\phi}(x, y) -\beta (\log \pi_{\theta}(y\mid x) - \log \pi_\text{ref}(y\mid x))}$, which is then maximized using Proximal Policy Optimization (PPO) \cite{schulman2017proximal}.

\section{Direct Preference Optimization}\label{sec:DPO}

Recognizing the complexities involved in scaling reinforcement learning approaches to large models during fine-tuning, we aim to formulate a streamlined policy optimization method that operates directly from preference data. Departing from conventional RLHF strategies, which entail reward learning followed by RL-based maximization, our method capitalizes on a specialized reward parameterization that permits explicit, closed-form derivation of the corresponding optimal policy—circumventing the need for an RL training loop. As elaborated in the next section, our central contribution is to establish an analytic transformation between the reward model and its induced policy, allowing us to translate losses over rewards into equivalent losses over policies. This reparameterization eliminates the necessity of independently fitting a reward function, yet continues to respect established models of human preference—such as the Bradley-Terry formulation

\textbf{Derivation of the DPO Objective.} We begin with the reinforcement learning (RL) objective provided in Eq.~\ref{eq:RL}, where the reward function $r$ is assumed to be arbitrary. As established in the literature~\citep{peters2007reinforcement, peng2019advantage, korbak2022reinforcement, go2023aligning}, the solution that optimally maximizes the KL-regularized reward (cf. Eq.~\ref{eq:RL}) assumes the form:

\begin{equation}\label{eq:op_policy}
    \pi_r(y\mid x) = \frac{1}{Z(x)}\pi_\text{ref}(y\mid x)\exp\left(\frac{1}{\beta}r(x, y)\right),
\end{equation}

with the normalization constant $Z(x) =\sum_{y}\pi_\text{ref}(y\mid x)\exp\left(\frac{1}{\beta}r(x, y)\right)$. Additional details can be found in Appendix \ref{app:derivation1}. Although one might utilize the MLE estimator $r_{\phi}$ to approximate the true reward $r^*$, computing the normalizer $Z(x)$ directly remains computationally intensive~\citep{korbak2022reinforcement, go2023aligning}, thus making direct use of Eq.~\ref{eq:op_policy} infeasible for practical implementation. To alleviate this, we can reformulate Eq.~\ref{eq:op_policy} by solving for the reward function in terms of its optimal policy $\pi_r$, reference policy $\pi_\text{ref}$, and the partition function $Z(\cdot)$. Explicitly, applying the logarithm to both sides and simplifying yields:

\begin{equation}\label{eq:main_eq}
    r(x,y) =\beta \log \frac{\pi_r(y\mid x)}{\pi_\text{ref}(y\mid x)} + \beta \log Z(x).
\end{equation}

This parameterization applies equally to the ground-truth reward $r^*$ and its associated optimal policy $\pi^*$. Importantly, the Bradley-Terry framework relies exclusively on reward differences between two candidates, i.e., ${p^*(y_1 \succ y_2 \mid x) = \sigma(r^*(x, y_1) - r^*(x, y_2))}$. Substituting the formulation from Eq.~\ref{eq:main_eq} into the Bradley-Terry preference structure (see Eq.~\ref{eq:bradley-terry}), we note that the partition function terms cancel, allowing the probability of observed preferences to be described solely using $\pi^*$ and $\pi_\text{ref}$. Thus, the optimal policy $\pi^*$ that solves RLHF under the Bradley-Terry likelihood meets the following criterion:

\begin{equation}\label{eq:objective}
    p^*(y_1\succ y_2 \mid x)=\frac{1}{1 + \exp\left(\beta \log \frac{\pi^*(y_2\mid x)}{\pi_\text{ref}(y_2\mid x)} - \beta \log \frac{\pi^*(y_1\mid x)}{\pi_\text{ref}(y_1\mid x)}\right)}
\end{equation}

A step-by-step derivation can be found in Appendix~\ref{app:derivation2}. While this is specific to the Bradley-Terry model, similar reasoning extends to the Plackett-Luce class~\citep{plackett1975analysis, luce2012individual}; see Appendix~\ref{app:plackett_luce_models} for details.

Since the probability of preference annotations is now formulated in terms of the optimal policy and not the reward, it is straightforward to propose a maximum likelihood estimation (MLE) loss for a parametric policy $\pi_\theta$. Drawing a parallel to the reward model objective (Eq.~\ref{eq:reward_model}), we derive the policy optimization loss as:

\begin{equation}\label{eq:optimum_model}
    \mathcal{L}_\text{DPO}(\pi_{\theta}; \pi_\text{ref}) = -\mathbb{E}_{(x, y_w, y_l)\sim \mathcal{D}}\left[\log \sigma \left(\beta \log \frac{\pi_{\theta}(y_w\mid x)}{\pi_\text{ref}(y_w\mid x)} - \beta \log \frac{\pi_{\theta}(y_l\mid x)}{\

In this framework, we approximate an implicit reward function via a distinct parameterization, where the corresponding optimal policy aligns with $\pi_\theta$. Additionally, because our methodology is tantamount to training a reparameterized Bradley-Terry model, it inherits certain desirable theoretical guarantees—such as consistency under appropriate assumptions regarding the distribution of preference data \cite{bong2022generalized}. We examine these theoretical aspects of DPO more thoroughly in Section~\ref{sec:theory}, drawing connections to related literature.

\textbf{Mechanics of the DPO update.} To gain insight into the DPO algorithm, it is instructive to investigate the gradient of its loss, $\mathcal{L}_\text{DPO}$. The derivative of $\mathcal{L}_\text{DPO}$ with respect to the model parameters $\theta$ is:

\begin{multline*}\label{eq:gradient}
    \nabla_\theta \mathcal{L}_\text{DPO}(\pi_\theta;\pi_\text{ref}) = \\ -\beta\mathbb{E}_{(x, y_w, y_l) \sim \mathcal{D}} \bigg[\underbrace{\sigma(\hat{r}_\theta(x, y_l) - \hat{r}_\theta (x, y_w))}_\text{places greater emphasis when reward ranking is inaccurate}\bigg[\underbrace{\nabla_\theta\log \pi(y_w \mid x)}_\text{promote $y_w$} - \underbrace{\nabla_\theta\log\pi(y_l \mid x)}_\text{demote $y_l$}\bigg]\bigg],
\end{multline*}

where $\hat{r}_\theta(x, y) = \beta \log \frac{\pi_\theta(y \mid x)}{\pi_\text{ref}(y \mid x)}$ denotes the reward implicitly modeled by the policy $\pi_\theta$ relative to a reference policy $\pi_\text{ref}$ (see Section~\ref{sec:theory} for more detail). Conceptually, this gradient encourages the model to assign higher probability to the preferred responses $y_w$ while reducing the likelihood assigned to the less-preferred responses $y_l$. Crucially, each sample is weighted according to the degree to which the implicit reward model $\hat{r}_\theta$ overrates the dispreferred responses, modulated by $\beta$. In other words, it measures the severity of the model's misordering of completions, with the magnitude determined by the KL regularization coefficient. Our experimental findings highlight the critical nature of this weighting scheme; omitting the weighting factor results in a straightforward variant of the algorithm that can lead to language model collapse (see Appendix Table~\ref{tab:unlikelihood_generations}).

\textbf{DPO summary.} The standard DPO workflow includes: 1) For each prompt $x$, generate completions $y_1, y_2 \sim \pi_\text{ref}(\cdot \mid x)$ and use human annotations to create an offline preference dataset $\mathcal{D} = \{x^{(i)}, y_w^{(i)}, y_l^{(i)}\}_{i=1}^N$; 2) update the language model $\pi_\theta$ by minimizing $\mathcal{L}_\text{DPO}$ with respect to the chosen $\pi_\text{ref}$, dataset $\mathcal{D}$, and the parameter $\beta$.  
In actual applications, practitioners often prefer to use existing, publicly available preference datasets rather than generating new completions and collecting human annotations from scratch. Because these datasets are generated from $\pi^\text{SFT}$, it is standard to set $\pi_\text{ref} = \pi^\text{SFT}$ if it exists. In cases where $\pi^\text{SFT}$ cannot be accessed, $\pi_\text{ref}$ is instead initialized by maximizing the likelihood over the preferred completions ${(x, y_w)}$, specifically, by setting ${\pi_\text{ref} = \argmax_{\pi}\mathbb{E}_{x, y_w \sim \mathcal{D}}\left[\log \pi(y_w \mid x)\right]}$. This strategy addresses the mismatch between the true but inaccessible reference distribution and the available $\pi_\text{ref}$ employed by DPO. Additional specifics regarding implementation and chosen hyperparameters are provided in Appendix~\ref{app:implementation}.

\section{Theoretical Analysis of DPO} This section provides a deeper interpretation of the DPO algorithm, outlines its theoretical basis, and discusses the strengths of DPO in comparison to the drawbacks of actor-critic approaches commonly used in RLHF, such as PPO~\cite{schulman2017proximal}.

\label{sec:theory}

\subsection{Your Language Model Is Secretly a Reward Model} DPO eliminates the need to learn an explicit reward model and for reinforcement learning steps, instead enabling direct policy learning through a single maximum likelihood objective. Importantly, the DPO loss (Eq. \ref{eq:main_eq}) is functionally equivalent to the Bradley-Terry model when rewards are parameterized as $r^*(x, y) = \beta \log\frac{\pi^*_\theta(y \mid x)}{\pi_\text{ref}(y \mid x)}$. The parameterized model $\pi_{\theta}$ is thus trained in a way analogous to fitting a reward model in Eq. \ref{eq:reward_model}, via an appropriate change of variables. The theory in this section clarifies the details of this reparameterization, demonstrates that it does not impose extra constraints on the space of learnable reward models, and ensures the possibility of exactly recovering the optimal policy. To this end, we begin by introducing an equivalence relationship for reward functions.

\begin{definition}
Two reward functions $r(x, y)$ and $r'(x, y)$ are considered equivalent if and only if there exists a function $f$ such that $r(x, y)-r'(x, y) = f(x)$.
\end{definition}

It is straightforward to verify that this definition satisfies the properties of an equivalence relation, thereby dividing the space of reward functions into equivalence classes. The following lemmas can then be established:

\begin{lemma}\label{lemma:same_prefrence} Within the Plackett-Luce preference framework, including the Bradley-Terry model as a special case, any two reward functions belonging to the same equivalence class induce identical preference distributions.
\end{lemma}

\begin{lemma}\label{lemma:same_policy}
    Any pair of reward functions in the same equivalence class will yield the same optimal policy in the associated constrained RL formulation.
\end{lemma}

The arguments supporting these lemmas are elementary, so their proofs are relegated to Appendix \ref{app:lemma1}. The first lemma restates a well-known issue of model under-specification within the Plackett-Luce family, as documented by \cite{plackett1975analysis}. This ambiguity necessitates the introduction of extra identifiability constraints in order to ensure meaningful MLE estimation, as discussed in Eq. \ref{eq:reward_model} and explored in \cite{bong2022generalized}. According to the second lemma, every reward function within the same equivalence class determines the same optimal policy; therefore, for purposes of policy recovery, it suffices to recover any representative reward function from the relevant class. In Appendix~\ref{app:thm1}, we formally establish the following result:

\begin{theorem}\label{thm:main}
    Provided mild conditions hold, any reward function compatible with Plackett-Luce or, specifically, Bradley-Terry models, can be reparameterized as ${r(x, y) = \beta \log \frac{\pi(y\mid x)}{\pi_\text{ref}(y\mid x)}}$ for a suitable model $\pi(y\mid x)$ and a fixed reference model $\pi_\text{ref}(y \mid x)$.
\end{theorem}

\begin{sproof}
    Let $r(x, y)$ be an arbitrary reward function associated with an optimal model $\pi_r(y \mid x)$ as given by Eq. \ref{eq:op_policy}. We aim to demonstrate that there exists a reward function in the equivalence class of $r$ expressible in the stated reparameterized form. Define the projection operator $f$ by  
\begin{equation}
    f(r; \pi_\text{ref}, \beta)(x, y) = r(x, y) - \beta\log\sum_{y}\pi_\text{ref}(y\mid x)\exp\left(\frac{1}{\beta}r(x, y)\right)
\end{equation}
This projection normalizes the reward by subtracting the log-partition function (dependent solely on $x$). Consequently, $f(r; \pi_\text{ref}, \beta)(x, y)$ resides within the equivalence class of $r(x, y)$. Substituting $r$ using the expression on the right-hand side of Eq.~\ref{eq:main_eq} (which universally applies to reward functions under this framework), it follows that $f(r; \pi_\text{ref}, \beta)(x, y) = \beta \log \frac{\pi_r(y\mid x)}{\pi_\text{ref}(y\mid x)}$. Therefore, this projection constructs a member of the equivalence class of $r$ in the required form, establishing that the reparameterization imposes no loss of generality for our class of reward models.
\end{sproof}

An alternative interpretation of Theorem~\ref{thm:main} is that it precisely identifies, for each equivalence class, the specific reward function chosen by the DPO reparameterization. Concretely, it selects the reward function that fulfills the condition:

\begin{equation}\label{eq:lag_p}
     \sum_{y}\underbrace{\pi_\text{ref}(y\mid x)\exp\left(\frac{1}{\beta}r(x, y)\right)}_{=\pi(y\mid x)\text{, using Thm.~\ref{thm:main} reparam.}} = 1,
\end{equation}

which ensures that $\pi(y\mid x)$ forms a legitimate probability distribution (namely, that all probabilities are non-negative and collectively sum to one). By referencing Eq.~\ref{eq:op_policy}, it becomes evident that Eq.~\ref{eq:lag_p} defines the partition function for the optimal policy corresponding to the reward $r(x, y)$. The central idea underpinning the DPO algorithm is to introduce constraints within the broadly defined and under-constrained Plackett-Luce family of preference models (including the Bradley-Terry model), thereby maintaining the set of representable reward models while ensuring that the optimal policy from Eq.~\ref{eq:op_policy} remains analytically manageable for any prompt $x$.

\subsection{Instability of Actor-Critic Algorithms} The proposed framework also facilitates the identification of sources of instability in commonly employed actor-critic methods for RLHF, such as PPO. Concentrating on the RL fine-tuning stage as described in Section~\ref{section:prelims}, we observe parallels with the control as inference formalism \cite{levine2018reinforcement} applied to the constrained reinforcement learning task presented in \ref{eq:RL}. Considering a parameterized policy $\pi_{\theta}(y\mid x)$, the goal is to minimize $\mathbb{D}_{\text{KL}}[\pi_{\theta}(y|x) \mid \mid \pi^*(y\mid x)]$, where $\pi^*$ represents the optimal policy induced by the reward $r_{\phi}(y, x)$ as given in Eq.~\ref{eq:optimum_model}. After rearranging terms, this results in the following objective for optimization:

\begin{equation}\label{eq:AC}
    \max_{\pi_{\theta}}\mathbb{E}_{\pi_{\theta}(y\mid x)}\bigg[\underbrace{r_{\phi}(x, y) -\beta\log\sum_{y}\pi_\text{ref}(y\mid x)\exp\left(\frac{1}{\beta}r_{\phi}(x, y)\right)}_{f(r_{\phi}, \pi_\text{ref}, \beta)} - \underbrace{\beta\log\frac{\pi_{\theta}(y\mid x)}{\pi_\text{ref}(y\mid x)}}_{\text{KL}}\bigg]
\end{equation}

The same objective function has been targeted in earlier studies 
\citep{ziegler2020finetuning, stiennon2022learning, bai2022training, ouyang2022training}, employing the DPO-equivalent reward within the reward family $r_{\phi}$. Here, the normalization component in $f(r_{\phi}, \pi_\text{ref}, \beta)$ may be viewed as the soft value function of the reference policy $\pi_\text{ref}$. Although this term does not alter the optimal policy, its exclusion can lead to high-variance gradients in policy optimization, thereby destabilizing training. To address this, one approach involves modeling the normalization with a learned value function, which introduces additional optimization challenges. Alternatively, earlier methods have resorted to reward normalization via a human baseline, effectively using a single sample as a Monte-Carlo estimator of the normalization factor. Distinctively, the DPO reparameterization leads to a reward function that eliminates the requirement for such baselines.

\section{Experiments}
Here, we provide an empirical assessment of DPO in preference-based policy training. To begin, in a controlled environment for text generation, we examine the efficiency with which DPO negotiates the balance between maximizing reward and minimizing KL-divergence relative to the reference policy—compared to widely adopted preference optimization algorithms such as PPO. Subsequently, we investigate DPO on larger-scale models and on more complex RLHF benchmarks, including tasks like summarization and dialogue. Our findings reveal that with minimal hyperparameter adjustment, DPO consistently matches or surpasses the performance of robust baselines such as RLHF with PPO, as well as outperforming the approach of selecting the best out of $N$ generated samples according to a learned reward. Prior to reporting these findings, we outline the experimental methodology; further technical details are provided in Appendix~\ref{app:exp_details}.

\textbf{Tasks.} We investigate three distinct open-ended text generation tasks in our experimental setup. Across all tasks, the algorithms aim to learn a policy using a dataset of preference pairs, represented as $\mathcal{D}=\bigl\{x^{(i)}, y_w^{(i)}, y_l^{(i)}\bigr\}_{i=1}^N$. In the case of \textbf{controlled sentiment generation}, the input $x$ consists of the beginning of a movie review drawn from the IMDb dataset \cite{maas-EtAl:2011:ACL-HLT2011}. Here, the objective is to generate an output $y$ that reflects a positive sentiment. To facilitate controlled analysis, for this particular task we synthesize pairs of generations and label preferences by applying a pre-trained sentiment classifier, ensuring that $p(\text{positive}\mid x,y_w)>p(\text{positive}\mid x,y_l)$. For SFT, we adapt GPT-2-large by fine-tuning it on the training set of IMDb reviews until performance plateaus (see App~\ref{app:sentiment_details} for more). For \textbf{summarization}, $x$ is a Reddit forum post, and the policy's task is to create a summary $y$ distilling the primary ideas in the original. In line with earlier studies, we utilize the Reddit TL;DR summarization dataset \citep{volske-etal-2017-tl}, supplemented with human preference annotations assembled by \citeauthor{stiennon2022learning}. The SFT baseline consists of a model that has been further trained on human-written summaries of forum posts,\footnote{\url{https://huggingface.co/CarperAI/openai_summarize_tldr_sft}} utilizing the TRLX \citep{leandro_von_werra_2023_7790115} toolkit for RLHF. Notably, the human preference annotations correspond to samples produced by a different, albeit comparably fine-tuned, SFT model as reported by \citeauthor{stiennon2022learning}. For the \textbf{single-turn dialogue} scenario, the input $x$ is a single user query, which may range from scientific inquiries to personal advice requests. The policy must respond with a reply $y$ that is both helpful and engaging, drawing upon the Anthropic Helpful and Harmless dialogue dataset \citep{bai2022training}, which contains 170,000 dialogues between a person and an AI assistant. At the end of each transcript, two candidate responses generated by a large (unspecified) language model are presented, along with an annotation identifying the preferred response from a human evaluator. Since a dedicated SFT model does not exist in this domain, we construct one by fine-tuning an available language model solely on the subset of completions marked as preferred.

\textbf{Evaluation.} Our experimental design employs two main evaluation strategies. To systematically investigate each algorithm’s capacity for maximizing rewards under constraints, we utilize a controlled sentiment generation framework. Here, we assess performance by charting the trade-off curve between achieved reward and KL-divergence from a baseline policy—made possible by access to an explicit, ground-truth reward signal (in this case, a sentiment classifier). In contrast, when moving to real-world scenarios where the ground-truth reward is inaccessible, our evaluation is based on the algorithm’s \textit{win rate} relative to a baseline policy, with GPT-4 serving as a surrogate for human judgments. Specifically, in the summarization and single-turn dialogue tasks, GPT-4 is used to judge the quality of summaries and the helpfulness of responses, respectively. For summarization, the reference summaries included in the test set establish the baseline, while in the dialogue task, the dataset’s preferred responses function as the baseline. Although prior work indicates that language models may outperform traditional metrics as evaluators \citep{Chen2023ExploringTU}, we corroborate the reliability of GPT-4 through a dedicated human-subjects study (see Sec.~\ref{sec:human-judgments}). This study reveals that GPT-4’s evaluations are closely aligned with those of human raters, and that agreement between humans and GPT-4 often matches or surpasses that between human annotators.

\textbf{Methods.} Beyond DPO, we examine a variety of established methods for training language models according to human-preference signals. As basic baselines, we include zero-shot prompting of \textbf{GPT-J} \citep{gpt-j} for the summarization task, and 2-shot prompting with \textbf{Pythia-2.8B} \citep{biderman2023pythia} for dialogue. Additionally, our evaluation spans the standard \textbf{SFT} approach and a \textbf{Preferred-FT} model, which involves fine-tuning on preferred outputs $y_w$—sourced from the SFT model in the context of sentiment and summarization, and from a standard LM for the single-turn dialogue task. Among alternative supervised strategies, we assess the \textbf{Unlikelihood} method \citep{welleck2019neural}, which encourages the model to maximize likelihood of $y_w$ while simultaneously suppressing the likelihood of $y_l$, using an optional coefficient $\alpha\in[0,1]$ applied to the unlikelihood component. We also investigate \textbf{PPO} \citep{schulman2017proximal} trained with a reward model inferred from preference annotations, as well as an oracle variant, \textbf{PPO-GT}, trained directly with the known ground-truth reward in the controlled sentiment experiments. For the latter, we implement both a standard version \cite{leandro_von_werra_2023_7790115} and a modified variant that includes reward normalization and more extensively tuned hyperparameters for better results; these enhancements are also applied to PPO models trained on learned rewards. Lastly, we include a \textbf{Best of $N$} comparison, where $N$ candidate outputs are sampled from the SFT model (or from Preferred-FT in dialogue), and the top-scoring one is selected based on the learned reward model. While this method achieves strong performance by separating the influence of the reward model from PPO optimization, it is computationally burdensome even for modest $N$ values, as it necessitates generating $N$ responses per query during testing.

\subsection{How well can DPO optimize the RLHF objective?}

Standard RLHF methodologies optimize a reward maximization objective that is constrained by a KL divergence term, promoting reward-seeking behaviors while preventing significant policy drift from a predefined reference policy. Consequently, algorithmic comparisons should jointly consider both the achieved reward and the magnitude of KL divergence; securing marginally greater reward at the cost of substantially increased KL divergence is not ideal. Figure~\ref{fig:frontier-tldr-main} illustrates the reward versus KL divergence tradeoff for multiple algorithms within the sentiment analysis context. For each algorithm, we perform several training runs, adjusting a hyperparameter governing the degree of policy conservatism per run (target KL $\in\{3,6,9,12\}$ for PPO, $\beta \in \{0.05,0.1,1,5\}$ for DPO, and $\alpha\in\{0.05,0.1,0.5,1\}$ for the unlikelihood approach; preferred-FT uses random seeds). The full sweep comprises 22 runs. After each increment of 100 training steps, up to convergence, we assess every policy on a test prompt set, recording both the average return with respect to the true reward metric and the mean sequence-level KL divergence\footnote{Namely, the aggregated KL across all timesteps in the sequence.} relative to the reference policy, $\text{KL}\left(\pi\mid \mid \pi_\text{ref}\right)$. The analysis reveals that DPO establishes the most favorable efficiency frontier, achieving maximal reward without incurring high KL divergence. This observation is particularly striking for several reasons. Primarily, both DPO and PPO target optimization of the same objective, yet DPO demonstrates considerably greater efficiency, with a reward-to-KL relationship that consistently surpasses that of PPO. Moreover, DPO outperforms PPO even when PPO is supplied with ground truth rewards (PPO-GT), yielding a superior reward/KL balance.

\subsection{Can DPO scale to real preference datasets?}
\label{sec:dpo-real-datasets}
We proceed to assess how DPO performs when fine-tuned on real-world preference data, specifically in the contexts of summarization and single-turn dialogue tasks. In the case of summarization, traditional automatic metrics like ROUGE are known to sometimes lack alignment with human judgments~\citep{stiennon2022learning}. Previous research has demonstrated that tuning language models with PPO using human preference data can yield more effective summaries. To evaluate this, we generate completions from various approaches on the test partition of the TL;DR summarization dataset and calculate their average win rate compared to the reference completions in the test data. Completions for each method are sampled using temperature values ranging between 0.0 and 1.0; the corresponding win rates are displayed in Figure~\ref{fig:frontier-tldr-main} (right). DPO, PPO, and Preferred-FT all use the same GPT-J SFT model as their starting point\footnote{\url{https://huggingface.co/CarperAI/openai_summarize_tldr_sft}}. Our results indicate that DPO attains an average win rate close to 61\% when the temperature is set to 0.0, which outperforms PPO’s peak performance of roughly 57\% at its optimal temperature, also 0.0. Furthermore, DPO achieves a greater highest win rate than the best-of-$N$ baseline method. It is important to highlight that no substantial tuning of the $\beta$ hyperparameter was undertaken for DPO; therefore, these results might not fully reflect the method’s capabilities. In addition, DPO shows significantly greater resilience to variations in sampling temperature than PPO, whose performance can drop to that of the unmodified GPT-J at higher temperature settings. Preferred-FT does not yield notable improvements over the base SFT model. In Section~\ref{sec:human-judgments}, we further present direct human comparison results between DPO and PPO, where DPO outputs at temperature 0.25 were preferred 58\% of the time compared to PPO outputs sampled at temperature 0.

For single-turn dialogue evaluation, we focus on a subset of the Anthropic HH dataset's test split \citep{bai2022training} containing single human-assistant exchanges. For GPT-4-based assessments, we utilize the preferred test completions as references to calculate each method's win rate. In the absence of a standard SFT baseline for this benchmark, we initialize with a pre-trained Pythia-2.8B model, train a reference model via Preferred-FT on selected completions to align output distributions, and subsequently apply DPO training. Comparisons are drawn against the highest-scoring completion from 128 Preferred-FT samples (having observed performance saturation at 128, as detailed in Appendix Figure~\ref{fig:best-of-n}), as well as against a two-shot prompted Pythia-2.8B base model; DPO achieves equivalent or superior outcomes at optimal temperatures relative to these baselines. Additionally, we assess an RLHF model optimized by PPO on the Anthropic HH dataset\footnote{\url{https://huggingface.co/reciprocate/ppo_hh_pythia-6B}} following guidance from a reputable implementation\footnote{\url{https://github.com/CarperAI/trlx/tree/main/examples/hh}}, but are unable to identify prompt or temperature configurations that surpass the original Pythia-2.8B's performance. Drawing on TL;DR findings and noting that both methods share an objective function, we interpret the Best of 128 metric as an approximate indicator of PPO performance. In summary, DPO uniquely combines computational efficiency with improvements over the Anthropic HH dataset's preferred completions, and delivers performance on par with or exceeding the much more resource-intensive Best of 128 strategy. Furthermore, Figure~\ref{fig:dialogue-main} demonstrates that DPO attains peak results with rapid convergence.

\subsection{Generalization to a new input distribution}

To more thoroughly investigate the robustness of PPO and DPO under distributional changes, we assess the policies trained on the Reddit TL;DR summarization task using samples from a different data domain: news articles contained in the test partition of the CNN/DailyMail dataset \citep{nallapati-etal-2016-abstractive}. For this evaluation, we adopt the optimal sampling temperatures identified on TL;DR (0 and 0.25). The findings, reported in Table~\ref{tab:ood}, show that GPT-4 win rates are computed with respect to the ground-truth summaries in each dataset, utilizing the same GPT-4 (C) prompt as in the Reddit TL;DR experiments, except with the phrase ``forum post'' modified to ``news article.'' On this alternate input distribution, DPO consistently surpasses the PPO policy by a notable margin. These results provide preliminary evidence that DPO can generalize as effectively as PPO, despite not relying on the extra unlabeled Reddit TL;DR prompts leveraged by PPO.

\subsection{Validating GPT-4 judgments with human judgments}
\label{sec:human-judgments}
To assess the trustworthiness of GPT-4’s evaluations, we conduct a human annotation study based on outcomes from the TL;DR summarization experiments, employing two variants of GPT-4 prompts. The \textbf{GPT-4 (S)} (simple) prompt inquires which summary better captures the key information from the post, while the \textbf{GPT-4 (C)} (concise) prompt also requests a preference for conciseness; this latter is included because GPT-4, when prompted with \textbf{GPT-4 (S)}, tends to prefer longer and more redundant summaries compared to human raters. Complete prompt wording can be found in Appendix~\ref{app:prompts}. Our evaluation involves three comparisons: the highest-performing method (DPO, temp. 0.25), the lowest-performing (PPO, temp. 1.0), and a mid-performing approach (SFT, temp. 0.25), each measured against greedily-sampled PPO (the best-performing temperature for PPO). This set is chosen to encompass a variety of summary qualities. Across both prompts, we observe that GPT-4's selections generally align with human preferences at similar rates as human annotators agree with one another, suggesting that GPT-4 provides a credible stand-in for human judgment (though for DPO and PPO-1 comparisons, multiple human ratings are gathered due to constraints on available raters). Broadly, the \textbf{GPT-4 (C)} prompt yields win rates that more closely track with human choices, and thus, this prompt is used as the default for the main analyses in Section~\ref{sec:dpo-real-datasets}. Additional details about the human annotation process, including the user interface shown to raters and the list of annotators, are provided in Appendix~\ref{app:human-study}.

\section{Discussion}
Preference-based learning presents an effective and extensible approach to developing language models that are both capable and well-aligned with human intent. In this work, we presented DPO, a straightforward method that enables language model training from preference data without employing reinforcement learning techniques. DPO circumvents the need to reformulate preference learning into a classical RL framework simply to leverage common RL tools. Instead, it establishes a direct correspondence between language model policies and reward functions, allowing the model to be trained directly to align with human preferences via a cross-entropy loss function—eliminating the reliance on reinforcement learning and preserving generality. DPO achieves performance on par with or superior to prevalent RLHF approaches, such as those based on PPO, with minimal hyperparameter tuning required. This simplicity and effectiveness markedly lower the barrier for leveraging human preference data to train language models.

\textbf{Limitations \& Future Work.} The findings presented prompt several avenues for further investigation. One open question concerns the extent to which DPO-trained policies generalize to novel distributions in comparison to models trained with explicit reward signals. Preliminary results indicate that generalization abilities of DPO policies are comparable to those derived from PPO, but a broader and more systematic analysis is necessary. Another area worth exploring is whether approaches such as self-labeling using DPO-trained policies can effectively utilize unlabeled data. In addition, further work is needed to understand how phenomena like reward over-optimization might appear in direct preference optimization, particularly in relation to the observed minor reduction in performance in Figure~\ref{fig:dialogue-main}-right, and whether this can be attributed to such effects. While our experiments included models up to 6B parameters, scaling DPO to much larger, state-of-the-art architectures remains an exciting research direction. With respect to evaluation methodologies, we observed that win rates measured by GPT-4 can vary with changes to the evaluation prompt, suggesting that optimizing prompt design for reliable judgments by automated systems is an important topic for future research. Lastly, the applicability of DPO extends beyond language modeling—future applications might include preference-based training of generative models in diverse domains and modalities.